% ============================================================================
% BAB III - IMPLEMENTASI DAN KONFIGURASI
% ============================================================================
\chapter{IMPLEMENTASI DAN KONFIGURASI}

Bab ini menjelaskan proses implementasi teknis meliputi konfigurasi
replikasi PostgreSQL, konfigurasi Pgpool-II sebagai
\textit{middleware}, serta integrasi aplikasi Laravel dengan
Pgpool-II. Setiap bagian disertai potongan konfigurasi dan kode
program beserta penjelasannya. Implementasi dilakukan secara bertahap
dimulai dari penyiapan PostgreSQL Master, kemudian Replica, lalu
Pgpool-II, dan terakhir integrasi dengan aplikasi. Seluruh konfigurasi
mengacu pada file aktual yang digunakan dalam proyek dan telah
diverifikasi berjalan pada lingkungan Docker Compose.

\section{Konfigurasi Replication}

Replikasi pada PostgreSQL dikonfigurasi menggunakan metode
\textbf{\textit{streaming replication}} asinkron. Master mengirimkan
data WAL (\textit{Write-Ahead Log}) ke setiap Replica secara
\textit{streaming} untuk menjaga agar data di Replica tetap
tersinkronisasi dengan Master. Konfigurasi melibatkan pengaturan pada
file \texttt{postgresql.conf} dan \texttt{pg\_hba.conf} di sisi
Master, serta \texttt{postgresql.conf} dan \texttt{entrypoint.sh} di
sisi Replica. Metode asinkron dipilih karena memberikan performa tulis
yang lebih tinggi dibandingkan replikasi sinkron, meskipun ada
kemungkinan \textit{replication lag} kecil antara Master dan Replica.
Untuk lingkungan pengembangan dan aplikasi \textit{e-commerce} yang
tidak memerlukan konsistensi \textit{real-time} absolut, pendekatan
ini sudah sangat memadai.

\subsection{Konfigurasi PostgreSQL Master}

\subsubsection{postgresql.conf (Master)}

File konfigurasi utama Master terletak di
\texttt{docker/pg-master/postgresql.conf}. Berikut adalah bagian
konfigurasi yang relevan dengan replikasi:

\begin{lstlisting}[caption={Konfigurasi Replikasi pada PostgreSQL Master
(docker/pg-master/postgresql.conf)}, label={lst:master-conf}]
# ===================== WAL (Write Ahead Log) =====================
wal_level = replica             # Level minimal untuk streaming replication
max_wal_senders = 10            # Maksimum koneksi replikasi bersamaan
wal_keep_size = 512             # Retensi WAL dalam MB
max_replication_slots = 10      # Maksimum replication slots

# ===================== Synchronous Replication =====================
synchronous_commit = on

# ===================== Logging =====================
log_statement = 'all'           # Mencatat semua query SQL (pembuktian)
log_line_prefix = '[%t] [%d] [%u] [%h] %p '
\end{lstlisting}

\textbf{Penjelasan konfigurasi:}
\begin{itemize}
    \item \texttt{wal\_level = replica}: Mengatur level WAL ke
    \texttt{replica} sebagai syarat minimal agar \textit{streaming
    replication} dapat berjalan. Level ini memastikan data yang cukup
    ditulis ke WAL untuk direplikasi, termasuk seluruh perubahan data
    dan informasi transaksi yang dibutuhkan Replica.
    \item \texttt{max\_wal\_senders = 10}: Menentukan jumlah maksimum
    koneksi \textit{WAL sender} secara bersamaan. Setiap Replica
    membutuhkan satu koneksi \textit{WAL sender} dari Master, sehingga
    nilai 10 cukup untuk mendukung hingga 10 Replica atau proses
    replikasi bersamaan lainnya.
    \item \texttt{wal\_keep\_size = 512}: Menentukan ukuran retensi
    file WAL dalam MB. Jika Replica tertinggal, WAL tetap disimpan
    hingga batas ukuran ini agar Replica dapat mengejar
    ketertinggalan tanpa perlu melakukan \texttt{pg\_basebackup}
    ulang dari awal.
    \item \texttt{max\_replication\_slots = 10}: Menentukan jumlah
    maksimum \textit{replication slot} untuk memastikan Master tidak
    menghapus WAL yang masih dibutuhkan oleh Replica. Slot replikasi
    memberikan jaminan bahwa WAL akan dipertahankan sampai semua
    Replica yang terdaftar telah mengonsumsinya.
    \item \texttt{log\_statement = 'all'}: Mencatat seluruh query SQL
    yang diterima Master. Digunakan sebagai bukti bahwa query
    \textit{write} diproses oleh Master dalam pengujian nanti.
\end{itemize}

\subsubsection{pg\_hba.conf (Master)}

File \texttt{docker/pg-master/pg\_hba.conf} mengatur autentikasi
koneksi yang masuk ke Master. File ini sangat kritis karena kesalahan
konfigurasi autentikasi akan menyebabkan Replica gagal terhubung dan
Pgpool-II tidak dapat melakukan \textit{health check}. Berikut adalah
konfigurasinya:

\begin{lstlisting}[caption={Konfigurasi pg\_hba.conf Master
(docker/pg-master/pg\_hba.conf)}, label={lst:master-hba}]
# Allow application user & replication user from Docker network
host    all             bdtuser     172.25.0.0/24    md5
host    all             postgres    172.25.0.0/24    md5
host    replication     repluser    172.25.0.0/24    md5
host    all             repluser    172.25.0.0/24    md5

# Allow Pgpool-II health check
host    all             repluser    172.25.0.20/32    trust
\end{lstlisting}

\textbf{Penjelasan konfigurasi:}
\begin{itemize}
    \item Baris \texttt{replication} mengizinkan user
    \texttt{repluser} melakukan koneksi replikasi dari seluruh
    \textit{container} dalam subnet 172.25.0.0/24. Metode
    \texttt{md5} memastikan bahwa password di-\textit{hash} selama
    proses autentikasi.
    \item Baris \textit{health check} mengizinkan Pgpool-II
    (172.25.0.20) melakukan pengecekan kesehatan tanpa
    \textit{password} (\texttt{trust}). Ini menyederhanakan koneksi
    periodik yang dilakukan Pgpool-II setiap 10 detik tanpa overhead
    autentikasi berulang.
    \item Autentikasi untuk \texttt{bdtuser} (user aplikasi)
    menggunakan metode \texttt{md5} untuk keamanan. Aplikasi Laravel
    akan mengirimkan password yang di-\textit{hash} MD5 saat
    menghubungi database melalui Pgpool-II.
\end{itemize}

\subsubsection{Pembuatan User Replikasi}

User replikasi dibuat melalui file inisialisasi
\texttt{docker/pg-master/init.sql} yang dijalankan otomatis saat
\textit{container} PostgreSQL Master pertama kali dijalankan. File ini
diekseskusi oleh mekanisme \texttt{/docker-entrypoint-initdb.d/}
bawaan image PostgreSQL Docker. Selain membuat user aplikasi
\texttt{bdtuser} dengan hak akses penuh, file ini juga membuat user
replikasi \texttt{repluser} yang akan digunakan oleh Replica untuk
melakukan \texttt{pg\_basebackup} dan \textit{streaming replication},
serta oleh Pgpool-II untuk \textit{health check}.

\begin{lstlisting}[language=SQL, caption={Inisialisasi User Replikasi
dan Aplikasi (docker/pg-master/init.sql)}, label={lst:init-sql}]
-- User aplikasi
CREATE ROLE bdtuser WITH LOGIN PASSWORD 'bdtpassword';
GRANT ALL PRIVILEGES ON DATABASE bdt_app TO bdtuser;

-- User replikasi (digunakan replica dan Pgpool health check)
CREATE ROLE repluser WITH LOGIN REPLICATION PASSWORD 'replpassword';
GRANT CONNECT ON DATABASE bdt_app TO repluser;
GRANT USAGE ON SCHEMA public TO repluser;
GRANT SELECT ON ALL TABLES IN SCHEMA public TO repluser;
\end{lstlisting}

User \texttt{repluser} memiliki hak akses \texttt{REPLICATION} yang
diperlukan untuk \textit{streaming replication} dan hak
\texttt{SELECT} pada tabel untuk keperluan \textit{health check}
Pgpool-II. Tanpa hak \texttt{SELECT} pada tabel, Pgpool-II tidak akan
dapat memverifikasi bahwa database berfungsi dengan benar saat
melakukan \textit{health check} periodik.

\subsection{Konfigurasi PostgreSQL Replica}

\subsubsection{postgresql.conf (Replica)}

File konfigurasi Replica terletak di
\texttt{docker/pg-replica/postgresql.conf}. Konfigurasi ini
memfokuskan pada mode \textit{standby} yang memungkinkan Replica
menerima koneksi \textit{read-only} sambil terus menerima dan
mengaplikasikan WAL dari Master. Parameter \texttt{hot\_standby}
adalah kunci yang membedakan \textit{warm standby} (tidak bisa
melayani query) dengan \textit{hot standby} (bisa melayani query
\texttt{SELECT}). Berikut adalah konfigurasinya:

\begin{lstlisting}[caption={Konfigurasi PostgreSQL Replica
(docker/pg-replica/postgresql.conf)}, label={lst:replica-conf}]
# ===================== Replication (Standby) =====================
wal_level = replica
hot_standby = on                  # Izinkan query read-only saat recovery
hot_standby_feedback = on         # Status replica dikirim ke master
max_wal_senders = 10
wal_receiver_timeout = 30s

# Sertakan auto-generated config dari pg_basebackup -R
# File ini berisi primary_conninfo untuk koneksi replikasi
include_if_exists = '/var/lib/postgresql/data/postgresql.auto.conf'

# ===================== Logging =====================
log_statement = 'all'             # Catat semua query (pembuktian)
log_line_prefix = '[%t] [%d] [%u] [%h] %p '
\end{lstlisting}

\textbf{Penjelasan konfigurasi:}
\begin{itemize}
    \item \texttt{hot\_standby = on}: Mengaktifkan mode
    \textit{hot standby} sehingga Replica dapat melayani query
    \texttt{SELECT} (read-only) selama proses \textit{recovery}
    maupun saat berjalan normal. Tanpa parameter ini, Replica hanya
    akan menjadi cadangan pasif yang tidak dapat melayani beban
    \textit{read}.
    \item \texttt{hot\_standby\_feedback = on}: Mengirimkan status
    Replica ke Master, mencegah Master menghapus WAL yang masih
    dibutuhkan Replica. Ini penting untuk menghindari konflik replikasi
    di mana Master menghapus WAL sebelum Replica sempat
    mengaplikasikannya.
    \item \texttt{primary\_conninfo}: Parameter ini tidak ditulis
    manual, melainkan dihasilkan otomatis oleh flag \texttt{-R} pada
    perintah \texttt{pg\_basebackup} di \texttt{entrypoint.sh}.
    Berisi informasi koneksi ke Master: \texttt{host=pg-master
    port=5432 user=repluser password=replpassword}.
    \item \texttt{log\_statement = 'all'}: Mencatat seluruh query SQL
    yang diterima Replica. Digunakan sebagai bukti bahwa query
    \textit{read} diproses oleh Replica dalam pengujian \textit{read-
    write splitting}.
\end{itemize}

\subsubsection{Script Inisialisasi Replica (entrypoint.sh)}

Setiap Replica diinisialisasi melalui script
\texttt{docker/pg-replica/entrypoint.sh} yang melakukan
\texttt{pg\_basebackup} dari Master. Script ini dirancang untuk
menangani dua skenario: inisialisasi pertama (melakukan
\texttt{pg\_basebackup}) dan restart (langsung menjalankan PostgreSQL
karena data sudah ada). Sebuah file \textit{sentinel} bernama
\texttt{replica\_initialized} digunakan untuk membedakan kedua
skenario tersebut, sehingga \texttt{pg\_basebackup} tidak diulang
setiap kali container direstart.

\begin{lstlisting}[language=bash, caption={Script Inisialisasi Replica
(docker/pg-replica/entrypoint.sh)}, label={lst:replica-entry}]
# Run base backup dari Master
export PGPASSWORD="replpassword"
pg_basebackup \
    -h pg-master \
    -p 5432 \
    -U repluser \
    -D "$DATA_DIR" \
    -Fp \
    -Xs \
    -P \
    -R
\end{lstlisting}

\textbf{Penjelasan flag \texttt{pg\_basebackup}:}
\begin{itemize}
    \item \texttt{-h pg-master}: Hostname Master sebagai sumber data.
    \item \texttt{-U repluser}: User dengan hak \texttt{REPLICATION}.
    \item \texttt{-D \$(DATA\_DIR)}: Direktori tujuan data Replica.
    \item \texttt{-Fp}: Format \textit{plain} (file biasa, bukan tar).
    \item \texttt{-Xs}: Menggunakan \textit{streaming} untuk mengambil
    WAL selama \textit{backup} berlangsung.
    \item \texttt{-P}: Menampilkan progres.
    \item \texttt{-R}: Membuat file \texttt{standby.signal} dan menulis
    \texttt{primary\_conninfo} ke \texttt{postgresql.auto.conf} secara
    otomatis.
\end{itemize}

\section{Konfigurasi Pgpool-II}

Pgpool-II bertindak sebagai \textit{middleware} yang berada di antara
aplikasi dan \textit{backend} PostgreSQL. Fungsinya mencakup tiga
aspek utama: \textit{connection pooling} (mengelola dan menggunakan
kembali koneksi ke backend), \textit{load balancing} (mendistribusikan
query SELECT ke beberapa Replica), dan \textit{read-write splitting}
(memisahkan query write ke Master dan query read ke Replica).
Konfigurasi Pgpool-II disimpan di \texttt{docker/pgpool/pgpool.conf}
dan beberapa parameter di-\textit{override} melalui \textit{environment
variables} di \texttt{docker-compose.yml} untuk fleksibilitas.

\subsection{Konfigurasi Backend}

Bagian konfigurasi \textit{backend} mendefinisikan seluruh node
PostgreSQL yang dikelola oleh Pgpool-II. Setiap \textit{backend}
memiliki empat parameter utama: hostname, port, weight (bobot load
balancing), dan flag. Urutan indeks backend (0, 1, 2) penting karena
backend 0 secara default dianggap sebagai Master dalam mode
\textit{streaming replication}. Berikut adalah konfigurasi backend:

\begin{lstlisting}[caption={Konfigurasi Backend Pgpool-II
(docker/pgpool/pgpool.conf)}, label={lst:pgpool-backend}]
# Backend 0: PostgreSQL Master
backend_hostname0 = 'pg-master'
backend_port0 = 5432
backend_weight0 = 0
backend_flag0 = 'ALLOW_TO_FAILOVER'

# Backend 1: PostgreSQL Replica 1
backend_hostname1 = 'pg-replica1'
backend_port1 = 5432
backend_weight1 = 1
backend_flag1 = 'ALLOW_TO_FAILOVER'

# Backend 2: PostgreSQL Replica 2
backend_hostname2 = 'pg-replica2'
backend_port2 = 5432
backend_weight2 = 1
backend_flag2 = 'ALLOW_TO_FAILOVER'
\end{lstlisting}

\textbf{Penjelasan parameter:}
\begin{itemize}
    \item \texttt{backend\_hostname}: Nama \textit{host} atau alamat IP
    dari setiap \textit{node} PostgreSQL. Pgpool-II menggunakan nama
    \textit{container} Docker sebagai \textit{hostname}, yang
    diresolusi oleh DNS internal Docker. Penggunaan hostname lebih
    direkomendasikan daripada IP statis untuk fleksibilitas.
    \item \texttt{backend\_port}: Port PostgreSQL pada \textit{node}
    tersebut (port internal container 5432). Karena semua komunikasi
    internal menggunakan port yang sama, Pgpool-II membedakan node
    berdasarkan hostname.
    \item \texttt{backend\_weight}: Bobot \textit{load balancing} untuk
    mendistribusikan query \texttt{SELECT}. Nilai \texttt{0} pada
    Master berarti Master \textbf{tidak} menerima query \textit{read}
    sama sekali, sehingga fokus sepenuhnya pada operasi
    \textit{write}. Nilai \texttt{1} pada Replica berarti
    masing-masing Replica mendapat porsi yang sama dalam distribusi
    query \texttt{SELECT}.
    \item \texttt{backend\_flag = 'ALLOW\_TO\_FAILOVER'}: Mengizinkan
    Pgpool-II untuk melepaskan (\textit{detach}) \textit{backend} yang
    gagal dan melanjutkan operasi dengan \textit{backend} yang tersisa.
    Tanpa flag ini, kegagalan satu \textit{backend} dapat menyebabkan
    Pgpool-II berhenti merespons.
\end{itemize}

\subsection{Konfigurasi Load Balancing dan Read-Write Splitting}

\begin{lstlisting}[caption={Konfigurasi Mode Pgpool-II
(docker/pgpool/pgpool.conf)}, label={lst:pgpool-mode}]
# Core Modes
load_balance_mode = on          # Aktifkan distribusi SELECT ke replica
master_slave_mode = on          # Aktifkan read-write splitting
master_slave_sub_mode = 'stream' # Deteksi replikasi via streaming

# Health Check
health_check_period = 10
health_check_timeout = 20
health_check_user = 'repluser'
health_check_password = 'replpassword'
health_check_database = 'bdt_app'

# Failover
failover_on_backend_error = on

# Logging (pembuktian read-write splitting)
log_connections = on
log_statement = on
log_per_node_statement = on     # Catat node mana yang memproses query
\end{lstlisting}

\textbf{Penjelasan parameter mode:}
\begin{itemize}
    \item \texttt{load\_balance\_mode = on}: Mengaktifkan mekanisme
    \textit{load balancing}. Query \texttt{SELECT} akan didistribusikan
    ke \textit{backend} dengan \texttt{weight} di atas 0 (yaitu Replica
    1 dan Replica 2). Tanpa mode ini, seluruh query akan diarahkan ke
    Master meskipun jenisnya \texttt{SELECT}.
    \item \texttt{master\_slave\_mode = on}: Mengaktifkan mekanisme
    \textit{read-write splitting}. Pgpool-II secara otomatis
    mengarahkan query \textit{write} ke Master (backend 0) dan query
    \textit{read} ke Replica berdasarkan analisis \textit{query
    pattern}. Mode ini adalah fondasi utama arsitektur yang dibangun.
    \item \texttt{master\_slave\_sub\_mode = 'stream'}: Menggunakan
    \textit{streaming replication} sebagai metode deteksi status
    Master/Replica. Pgpool-II memeriksa \texttt{pg\_stat\_replication}
    untuk menentukan node mana yang berperan sebagai primary dan mana
    yang standby.
    \item \texttt{log\_per\_node\_statement = on}: Mencatat
    \textit{node} backend mana yang memproses setiap \textit{statement}
    SQL. Ini sangat penting sebagai bukti bahwa \textit{read-write
    splitting} berjalan. Setiap query akan tercatat dengan label
    \texttt{DB node N} yang menunjukkan indeks backend.
    \item \texttt{health\_check\_*}: Pgpool-II secara periodik (setiap
    10 detik) memeriksa kesehatan setiap \textit{backend} menggunakan
    user \texttt{repluser}. Jika \textit{backend} gagal setelah 3 kali
    percobaan (\texttt{health\_check\_max\_retries}), \textit{backend}
    tersebut di-\textit{detach} dari pool.
    \item \texttt{failover\_on\_backend\_error = on}: Jika terjadi
    kesalahan pada \textit{backend}, Pgpool-II akan memicu
    \textit{failover} dan mengalihkan \textit{traffic} ke
    \textit{backend} yang masih sehat. Fitur ini diuji pada skenario B
    di Bab IV.
\end{itemize}

\subsection{Konfigurasi Autentikasi Client (pool\_hba.conf)}

\begin{lstlisting}[caption={Konfigurasi pool\_hba.conf
(docker/pgpool/pool\_hba.conf)}, label={lst:pool-hba}]
# TYPE  DATABASE    USER        ADDRESS          METHOD
local   all         all                         trust
host    all         all         127.0.0.1/32    trust
host    all         bdtuser     172.25.0.0/24    md5
host    all         all         0.0.0.0/0        md5
\end{lstlisting}

File \texttt{pool\_hba.conf} mengatur autentikasi \textit{client} yang
terhubung ke Pgpool-II. Mirip dengan \texttt{pg\_hba.conf} pada
PostgreSQL, file ini mendefinisikan aturan akses berdasarkan tipe
koneksi, database, user, alamat IP, dan metode autentikasi. User
\texttt{bdtuser} dari jaringan Docker diizinkan dengan metode
\texttt{md5} untuk keamanan. Aturan terakhir dengan \texttt{0.0.0.0/0}
memungkinkan koneksi dari luar jaringan Docker dengan metode
\texttt{md5}.

\section{Integrasi Aplikasi}

Aplikasi Laravel dikonfigurasi untuk terhubung ke
\textbf{Pgpool-II}, bukan langsung ke PostgreSQL Master atau Replica.
Ini memastikan seluruh mekanisme \textit{read-write splitting} dan
\textit{load balancing} berjalan secara transparan tanpa perlu
modifikasi kode aplikasi. Dari sudut pandang aplikasi, Pgpool-II
tampak seperti server PostgreSQL biasa yang berjalan pada host
\texttt{pgpool} dan port \texttt{5432}. Aplikasi tidak perlu
mengetahui bahwa di belakang Pgpool-II terdapat tiga node PostgreSQL
dengan peran berbeda.

\subsection{Konfigurasi Environment}

File \texttt{PWEB-B1/.env} (dan \texttt{.env.bdt} untuk Docker)
mengatur koneksi database. File \texttt{.env} digunakan saat
pengembangan lokal, sementara \texttt{.env.bdt} digunakan saat
aplikasi berjalan di dalam container Docker. Kedua file menetapkan
nilai \texttt{DB\_HOST=pgpool} yang merupakan hostname Pgpool-II dalam
jaringan Docker.

\begin{lstlisting}[caption={Konfigurasi Environment Aplikasi
(PWEB-B1/.env)}, label={lst:env}]
DB_CONNECTION=pgsql
DB_HOST=pgpool            # Koneksi ke Pgpool-II, bukan langsung ke Master
DB_PORT=5432
DB_DATABASE=bdt_app
DB_USERNAME=bdtuser
DB_PASSWORD=bdtpassword
\end{lstlisting}

\textbf{Catatan penting:} Nilai \texttt{DB\_HOST=pgpool} menunjukkan
bahwa aplikasi terhubung ke Pgpool-II (172.25.0.20), bukan ke
PostgreSQL Master (172.25.0.10) atau Replica (172.25.0.11/12).
Pgpool-II-lah yang akan meneruskan query ke \textit{backend} yang
sesuai berdasarkan aturan \textit{read-write splitting} yang telah
dikonfigurasi.

\subsection{Konfigurasi Datasource Laravel}

File \texttt{PWEB-B1/config/database.php} mendefinisikan koneksi
PostgreSQL dalam array konfigurasi Laravel. Konfigurasi ini membaca
nilai dari \textit{environment variable} yang telah diatur di file
\texttt{.env}, sehingga nilai koneksi dapat berubah antara lingkungan
pengembangan dan produksi tanpa mengubah kode.

\begin{lstlisting}[language=PHP, caption={Konfigurasi Datasource
Laravel (PWEB-B1/config/database.php)}, label={lst:datasource}]
'pgsql' => [
    'driver' => 'pgsql',
    'url' => env('DATABASE_URL'),
    'host' => env('DB_HOST', '127.0.0.1'),
    'port' => env('DB_PORT', '5432'),
    'database' => env('DB_DATABASE', 'forge'),
    'username' => env('DB_USERNAME', 'forge'),
    'password' => env('DB_PASSWORD', ''),
    'charset' => 'utf8',
    'prefix' => '',
    'prefix_indexes' => true,
    'search_path' => 'public',
    'sslmode' => 'prefer',
],
\end{lstlisting}

Konfigurasi datasource membaca nilai dari \textit{environment
variable} yang telah diatur di file \texttt{.env}.
\texttt{DB\_HOST} diambil dari \texttt{env('DB\_HOST')} yang bernilai
\texttt{pgpool}. Dengan mekanisme \texttt{env()} ini, koneksi database
dapat dengan mudah dialihkan ke server lain hanya dengan mengubah file
\texttt{.env} tanpa menyentuh kode sumber aplikasi.

\subsection{Koneksi di docker-compose.yml}

Pada level orkestrasi, \textit{environment variable} aplikasi diatur
di \texttt{docker-compose.yml} pada bagian \texttt{environment} dari
service \texttt{app}. Ini memastikan bahwa saat container aplikasi
dijalankan melalui Docker Compose, seluruh variabel koneksi database
sudah terkonfigurasi dengan benar.

\begin{lstlisting}[caption={Environment Variable Aplikasi di
docker-compose.yml}, label={lst:compose-env}]
app:
  environment:
    DB_CONNECTION: pgsql
    DB_HOST: pgpool               # Aplikasi KONEK ke PGPool
    DB_PORT: "5432"
    DB_DATABASE: bdt_app
    DB_USERNAME: bdtuser
    DB_PASSWORD: bdtpassword
\end{lstlisting}

Dengan konfigurasi di atas, seluruh koneksi database dari aplikasi
Laravel mengalir melalui Pgpool-II yang kemudian secara transparan
mengarahkan query \textit{write} ke Master dan query \textit{read} ke
Replica berdasarkan mekanisme \textit{read-write splitting}. Aplikasi
tetap menggunakan sintaks SQL standar tanpa perlu memodifikasi kode
untuk menentukan target server, karena seluruh logika \textit{routing}
ditangani sepenuhnya oleh Pgpool-II.
